% TOP_PROBLEM: {"problemId": "problem.anosov-diffeomorphism-classification-conjecture", "canonicalTitle": "Anosov Diffeomorphism Classification Conjecture", "releaseRank": 255, "primaryDomain": "probability_ergodic_dynamics", "primaryDomainLabel": "Probability, ergodic theory & dynamics", "displayStatus": "open", "releaseStatus": "open"}
% !TeX program = lualatex
% Definition and sources only; this file is not an LLM solution attempt.
\documentclass[11pt]{article}
\usepackage[margin=25mm]{geometry}
\usepackage{fontspec}
\setmainfont{Latin Modern Roman}
\usepackage{amsmath,amssymb,mathtools}
\usepackage{unicode-math}
\setmathfont{Latin Modern Math}
\usepackage{xurl,hyperref}
\hypersetup{colorlinks=true,urlcolor=blue}
\setcounter{secnumdepth}{0}
\setlength{\parindent}{0pt}
\setlength{\parskip}{5pt}
\setlength{\emergencystretch}{3em}
\begin{document}
\section*{\texorpdfstring{Anosov Diffeomorphism Classification Conjecture}{Anosov Diffeomorphism Classification Conjecture}}\label{255}
\subsection{Definitions and mathematical statement}
An Anosov diffeomorphism $f:M\to M$ of a closed smooth manifold has an invariant splitting $TM=E^s\oplus E^u$ and constants $C>0$, $0<\lambda<1$ with
$\|Df^nv^s\|\le C\lambda^n\|v^s\|$ and
$\|Df^{-n}v^u\|\le C\lambda^n\|v^u\|$ for $n\ge0$.
An infranilmanifold is a compact quotient of a simply connected nilpotent Lie group by a torsion-free discrete affine group with finite holonomy. The conjecture asks for a homeomorphism $h$ and a hyperbolic infranil automorphism $A$, with no eigenvalue of its inducing differential on the unit circle, such that
\[
 h\circ f=A\circ h.
\]
The precise infranil automorphism convention follows the source; topological conjugacy, not differentiable conjugacy, is requested.
\par\smallskip\textit{Formulation and concept references:} \href{https://aimath.org/WWN/measrigid/measrigid.pdf}{[S1]}.

\subsection{Short English statement}
Is every uniformly hyperbolic diffeomorphism topologically equivalent to an algebraic hyperbolic automorphism on an infranilmanifold?
\subsection{Sources}
\begin{itemize}
\item[S1]Open Problems from the Workshop 'Emerging Applications of Measure Rigidity' (AIM, 2004), Section 2.1, Conjecture 3.
\url{https://aimath.org/WWN/measrigid/measrigid.pdf}.
\textit{Formulation source.}
\end{itemize}
\end{document}
